\section{Experimental setup}
\label{sec:setup}

\subsection{Model ladders}
\label{sec:setup:models}

A scale effect means something only within a family trained with one recipe, so we evaluate ladders rather than a collection of popular models (Table~\ref{tab:ladders}). Every ladder stops at 8B parameters, the largest size whose 16-bit weights, about 16~GB, leave room for everything else in the laptop's 32~GB (Section~\ref{sec:setup:hw}). Qwen2.5 \citep{qwen2024qwen25} is the primary ladder, with four dense sizes from 0.5B to 7B released together. Qwen3 \citep{yang2025qwen3} adds four sizes from 0.6B to 8B whose post-trained checkpoints switch thinking on and off inside the same weights; we run them with thinking off everywhere except in the reasoning ablation. SmolLM2 \citep{allal2025smollm2} and Gemma 3 \citep{gemma2025gemma3} extend the low end to 135M and 270M parameters, with three sizes each. At every size we run both the base (pre-trained) and the instruct (post-trained) checkpoint, 28 checkpoints in all. Two fully open ladders serve as controls in a separate configuration that we run if time allows: Pythia \citep{biderman2023pythia}, six base-only sizes from 70M to 2.8B trained on the same data in the same order, and OLMo 2 \citep{olmo2025olmo2}, with open data and base and instruct checkpoints at 1B and 7B. The larger rungs, Qwen2.5 from 14B to 72B, Qwen3 at 14B and 32B, Gemma 3 at 12B and 27B, Pythia at 6.9B and 12B and OLMo 2 at 13B and 32B, are left to future work (Section~\ref{sec:future}).

We leave out ladders that would put different kinds of model on one axis. The Llama 3 herd \citep{grattafiori2024llama3} covers 8B, 70B and 405B, and the smaller Llama sizes come from later releases; the DeepSeek-R1 distillations sit on base models from two families \citep{deepseekai2025r1}; and mixture-of-experts models would force a choice between active and total parameters. Newer dense ladders, such as Olmo 3 \citep{olmo2025olmo3}, can be added through the model catalog once their repository identifiers are verified.

\begin{table}[t]
\centering\small
\caption{Model ladders in this study, all run with 16-bit weights on one laptop. ``Main'' is the main sweep; ``LoRA'' and ``free reply'' are the two configurations that complete it, and ``controls'' the optional control configuration (Section~\ref{sec:setup:factors}). Rungs above 8B are discussed in Section~\ref{sec:future}, and Appendix~\ref{app:ladders} gives the repository identifiers.}
\label{tab:ladders}
\begin{tabularx}{\linewidth}{@{}l X l X X@{}}
\toprule
Family & Sizes & Checkpoints & Configurations & License\\
\midrule
Qwen2.5 & 0.5B, 1.5B, 3B, 7B & base, instruct & main; LoRA up to 3B; free reply (instruct, every size) & Apache-2.0; 3B under the Qwen Research License\\
Qwen3 & 0.6B, 1.7B, 4B, 8B & base, post-trained & main, thinking off; LoRA up to 4B; free reply (post-trained, every size); thinking on in the reasoning ablation & Apache-2.0\\
SmolLM2 & 135M, 360M, 1.7B & base, instruct & main; LoRA; free reply (instruct, 360M and 1.7B) & Apache-2.0\\
Gemma 3 & 270M, 1B, 4B & pre-trained, instruct & main; LoRA; free reply (instruct, 1B and 4B) & Gemma Terms of Use\\
Pythia & 70M, 160M, 410M, 1B, 1.4B, 2.8B & base only (deduplicated) & controls & Apache-2.0\\
OLMo 2 & 1B, 7B & base, instruct & controls & Apache-2.0\\
\bottomrule
\end{tabularx}
\end{table}

Every checkpoint runs with 16-bit weights, bfloat16 where the Metal backend supports it and float16 otherwise, and every table marks the precision of each row. No rung is quantized, so the top of each ladder does not mix scale with quantization, as it would if the largest models were loaded with 4-bit weights; the price is that the ladders end at 8B. The parameter axis is the total parameter count from the model catalog. Embeddings take a large share of the smallest models, about a quarter of Qwen2.5-0.5B's 0.49B parameters according to its model card, and more in models with large vocabularies, so we repeat the slope fits with non-embedding counts where they are published.

\subsection{Factors and levels}
\label{sec:setup:factors}

Table~\ref{tab:factors} lists the factors we vary and the configuration that varies each. The main sweep, \texttt{configs/mac\_main.yaml}, crosses the 28 checkpoints with four methods on the seven tasks: generation with a 16-token budget, likelihood scoring, the probe at the middle layer trained on 100 episodes, and the feature probe, which needs no model and runs once per task. Two configurations complete it where the full cross would not fit on the laptop. \texttt{configs/mac\_lora.yaml} runs LoRA behavior cloning with 40 training episodes on the 24 checkpoints of at most 4B, and \texttt{configs/mac\_free\_reply.yaml} runs the free reply, with up to 96 new tokens, on twelve instruct or post-trained checkpoints (SmolLM2 at 360M and 1.7B, the four Qwen2.5 and the four Qwen3 sizes, and Gemma 3 at 1B and 4B) on four tasks, loan approval, support triage, gridworld and tic-tac-toe, because free replies are slow to generate. The control configuration, \texttt{configs/mac\_controls.yaml}, runs Pythia and OLMo 2 with the four methods of the main sweep.

Five ablation blocks under \texttt{configs/ablations/} vary one factor at a time, on Qwen2.5-Instruct at 1.5B, 3B and 7B unless stated otherwise. The prompting block (\texttt{prompting.yaml}) adds three demonstrations, letter labels, sampling at temperature 0.7, a contextual-bandit history window of 3 instead of 10, length-normalized scoring and scoring with demonstrations, on all seven tasks. The reasoning block (\texttt{reasoning.yaml}) gives chain-of-thought a budget of 256 new tokens on the three Qwen2.5 models and compares Qwen3 at 1.7B, 4B and 8B with thinking on, also with 256 tokens, against thinking off with 16; generating 256 tokens per decision is slow on the laptop, so this block runs on gridworld, tic-tac-toe and support triage only. The adaptation block (\texttt{adaptation.yaml}) moves to the three smallest Qwen2.5-Instruct sizes, 0.5B, 1.5B and 3B, and trains probes on 10, 40 and 100 episodes, probes at the last layer and at layer 4, the feature probe, and LoRA on 10 and 40 episodes, and on 100 at 0.5B and 1.5B, on all seven tasks. The task-structure block (\texttt{tasks.yaml}) runs thirteen task configurations, among them the shifted loan variants, with generation, scoring, the probe and the feature probe, and with LoRA on 40 episodes at 1.5B and 3B.

The quantization block (\texttt{quantisation.yaml}) takes a different route, because the usual one is closed on Apple silicon. An earlier draft of this design loaded 4-bit NF4 and 8-bit weights through bitsandbytes \citep{dettmers2023qlora,dettmers2022llmint8}, which needs CUDA. We compare instead the 16-bit weights of the local backend with two GGUF quantizations of the same checkpoints: Q8\_0, with 8-bit weights, and Q4\_K\_M, a 4-bit format that keeps some tensors at higher precision. llama.cpp serves the GGUF files with its Metal backend on the same laptop, and the harness reaches them through its OpenAI-compatible backend. That interface returns only the top log-probabilities of the first generated token, so neither the probe nor LoRA can be applied, and exact scoring becomes first-token scoring. The block therefore runs generation and first-token scoring with letter labels, at all three sizes and on all seven tasks; the 16-bit reference runs the same two conditions locally, where scoring single-token letter labels is the same as scoring their first token. H5 is restated for these formats and these two methods (Section~\ref{sec:res:h5}).

\begin{table}[t]
\centering\small
\caption{Factors, their levels in the configuration files, and the hypotheses they test. ``Main'', ``LoRA'', ``free reply'' and ``controls'' are \texttt{configs/mac\_main.yaml}, \texttt{mac\_lora.yaml}, \texttt{mac\_free\_reply.yaml} and \texttt{mac\_controls.yaml}; ``ablations'' are the five blocks under \texttt{configs/ablations/}.}
\label{tab:factors}
\begin{tabularx}{\linewidth}{@{}l X l l@{}}
\toprule
Factor & Levels & Configuration & Tests\\
\midrule
Parameter count & SmolLM2 135M to 1.7B; Qwen2.5 0.5B to 7B; Qwen3 0.6B to 8B; Gemma 3 270M to 4B; controls: Pythia 70M to 2.8B, OLMo 2 1B and 7B & main, controls & H1\\
Instruction tuning & base vs.\ instruct at each of the 14 sizes of the four ladders; OLMo 2 at 1B and 7B & main, LoRA, controls & H2\\
Conversion method & \texttt{prompt\_generate}, \texttt{prompt\_score}, \texttt{probe} and \texttt{feature\_probe} on all 28 checkpoints; \texttt{lora\_sft} on the 24 of at most 4B; \texttt{prompt\_free} on 12 & main, LoRA, free reply & H3, H4\\
Quantization & 16-bit vs.\ GGUF Q8\_0 and Q4\_K\_M at 1.5B, 3B and 7B, under generation and first-token letter scoring & ablations & H5\\
History format & bandit history as a per-arm summary vs.\ the raw list & ablations & H6\\
Reasoning budget & direct answer (16 tokens) vs.\ chain-of-thought (256) on Qwen2.5-Instruct; Qwen3 thinking off (16) vs.\ on (256) at 1.7B, 4B and 8B & ablations & H7\\
Prompt details & names vs.\ letters; 0 vs.\ 3 demonstrations; temperature 0 vs.\ 0.7; history window 10 vs.\ 3; raw vs.\ length-normalized scores; scoring with demonstrations & ablations & secondary\\
Adaptation data & probe on 10, 40 or 100 episodes, at the middle layer, the last layer or layer 4; LoRA on 10 or 40 episodes, and 100 at 0.5B and 1.5B & ablations & H3\\
Task difficulty & 5 vs.\ 10 arms; arm gap 0.2 vs.\ 0.5; grid 6 vs.\ 9; random vs.\ minimax opponent; 20 vs.\ 40 tickets; loan shift none, covariate or sign flip & ablations & H1, H3\\
\bottomrule
\end{tabularx}
\end{table}

\subsection{Seeds and data splits}
\label{sec:setup:seeds}

The main sweep, the LoRA and control configurations and the decision-model comparison evaluate seeds 0 to 99 for every task except tic-tac-toe, which uses seeds 0 to 299 because a game's outcome is $-1$, 0 or $+1$; blackjack keeps seeds 0 to 99 but plays twenty hands per episode for the same reason. The ablations and the free-reply configuration use the first half of each range: seeds 0 to 49, 150 tic-tac-toe games, and 50 blackjack episodes of twenty hands. Probe and LoRA training states and few-shot demonstrations come only from seeds 100000 to 100099, LoRA using the first 40 of them, and the harness rejects any configuration in which training and evaluation seeds overlap. Every model and method sees the same evaluation seeds, so all comparisons within a task configuration are paired, and an ablation cell is compared with a main-sweep cell over the seeds they share. The pilot (\texttt{configs/pilot.yaml}), in which prompts are debugged, evaluates six small checkpoints on seeds 5000 to 5019 with training states from the same training range, so no confirmatory instance is seen while prompts are being edited.

\subsection{Hardware and software}
\label{sec:setup:hw}

Every run in this study is made on one laptop, an Apple M2 Max with 32~GB of unified memory, through PyTorch's Metal (MPS) backend. The choice is deliberate: a machine like this is within reach of many researchers and practitioners, and a study that runs end to end on it can be checked by someone without a cluster or a cloud budget. Models run through the Transformers library \citep{wolf2020transformers} with 16-bit weights, bfloat16 where Metal supports it and float16 otherwise; LoRA runs through PEFT and probes through scikit-learn \citep{pedregosa2011scikit}. Nothing is loaded with 4-bit or 8-bit weights through bitsandbytes, which is not available on Apple silicon. Two parts of the study run outside the local backend, on the same laptop: the quantization block, whose GGUF files llama.cpp serves (Section~\ref{sec:setup:factors}), and the decision models, which run in their own servers and packages (Section~\ref{sec:setup:decision}).

The machine sets the boundaries of the design. Exact likelihood scoring, hidden-state probes and LoRA need the weights inside the local backend, and 16-bit weights take about two bytes per parameter, about 16~GB at 8B. With activations, the operating system and the harness sharing the same unified memory, these methods stop at about 8B parameters. LoRA stops earlier, at 4B, because of training time rather than memory, and even there it is trained on 40 episodes capped at 3000 states. Section~\ref{sec:future} lists what these limits leave out.

The run manifest records the platform string and the Python, PyTorch, Transformers, PEFT, scikit-learn and NumPy versions, so every result directory carries its own software record: \result{[macOS, PyTorch and Transformers versions from run\_manifest.json, and the llama.cpp build of the quantization block]}. Greedy generation, scoring and the probe involve no sampling, so reruns on the same laptop and library versions should agree up to the numerical nondeterminism of the Metal kernels, which we check by repeating \result{[number of cells repeated, and the largest difference in \norm{}]}; across hardware we expect agreement within confidence intervals, not identical values.

Latencies are wall-clock times per decision on this laptop. PyTorch's Metal backend is slower than runtimes written for Apple silicon, such as MLX or llama.cpp, so the absolute numbers understate what the same machine can do. Our latency comparisons are made within one machine and one runtime, and their ratios carry over to other settings better than their values. The quantized models in llama.cpp and the decision models in their own servers and packages share the machine but not the runtime, and their latencies are reported separately.

\subsection{Statistical protocol}
\label{sec:setup:stats}

The unit of analysis is the episode. The primary outcome is the normalized score \norm{} of Equation~\eqref{eq:norm}; the secondary outcomes are the illegal-action rate, oracle agreement, the task metrics, and latency and tokens per decision. For each cell we report the mean and the interquartile mean of \norm{} over episodes and a 95\% percentile bootstrap interval from 2000 resamples of episodes \citep{efron1993bootstrap,agarwal2021precipice}.

Two cells of the same task configuration are compared through the paired differences $\delta_\sigma=\norm(\pi_A,\sigma)-\norm(\pi_B,\sigma)$ over their shared seeds. The test statistic is $|\bar\delta|$, its null distribution comes from 5000 random sign flips of the $\delta_\sigma$, and the two-sided $p$-value is $(1+\#\{k: |\bar\delta^{(k)}|\ge|\bar\delta|\})/5001$. Effect sizes are paired Cohen's $d=\bar\delta/\operatorname{sd}(\delta)$. Within each hypothesis, the tests named in its pre-registered criterion form one family and are corrected with Holm's step-down procedure \citep{holm1979simple}; \texttt{deadeye compare} takes one family of pairs per call and applies the correction within it.

Scale is summarized by the least-squares fit
\begin{equation}
\label{eq:slope}
\bar{s}_c = a + b\,\log_{10} N_c
\end{equation}
over the cells $c$ of one ladder, method and task, where $\bar s_c$ is the cell mean of \norm{} and $N_c$ the parameter count; $b$ is the change in normalized score per tenfold increase in parameters. Each of the four ladders has three or four rungs between 135M and 8B, and \texttt{deadeye report} fits a slope only where there are at least three. Confidence intervals for $b$, for differences of slopes (H1b) and for ratios of slopes (H4b) come from refitting after resampling episodes within every cell, 1000 times; \texttt{deadeye report} writes the fits to \texttt{slopes.csv}, including a fit on the cross-task mean score per model, and \texttt{deadeye.report.slope\_ratio} forms the ratios. The fit is made separately for the base and the instruct checkpoints of a ladder (\texttt{slopes.csv} carries a column for the checkpoint type), so H1 is tested on both; a ladder enters the fit only where that type has at least three rungs.

The smallest effect of interest is $\Delta=0.10$ in \norm{}. For a paired comparison in which the per-episode differences have standard deviation $\sigma_\delta$, the number of episodes needed for power $1-\beta$ at level $\alpha$ is
\begin{equation}
\label{eq:power}
n = \Bigl\lceil \bigl((z_{1-\alpha/2}+z_{1-\beta})\,\sigma_\delta/\Delta\bigr)^2 \Bigr\rceil,
\end{equation}
which \texttt{deadeye.metrics.episodes\_needed} implements. With $\alpha=0.05$ and power 0.8 this gives 32 episodes at $\sigma_\delta=0.2$ and 126 at $\sigma_\delta=0.4$, so the 100 episodes of the main sweep suffice when $\sigma_\delta$ is at most about 0.36. The ablations and the free-reply configuration evaluate 50 episodes per cell to fit the laptop's time budget, which at the same $\sigma_\delta$ of 0.36 detects differences of about 0.14 rather than 0.10, and about 0.20 at $\sigma_\delta=0.5$. We report the smallest effect each block can detect next to its results and do not read a non-significant ablation contrast as the absence of an effect. The pilot estimates of $\sigma_\delta$ are \result{[per-task SD of paired differences from the pilot]}. Tasks whose anchors lie close together have a much larger per-episode spread of \norm{}. For blackjack the per-hand standard deviation of the return is close to 1, so the standard deviation of \norm{} is roughly the reciprocal of the anchor span \result{[span $\bar R_{\mathrm{orc}}-\bar R_{\mathrm{rand}}$ on blackjack]}, which is why the sweep bundles twenty hands into each of its 100 blackjack episodes and runs 300 tic-tac-toe games: the standard deviation of an episode mean falls with the square root of the number of hands. The minimum detectable difference under this design is \result{[per task, at power 0.8, for 100 and for 50 episodes]}, and where it exceeds 0.10 we report it next to the estimate rather than reading a non-significant result as a null.

A cell is excluded only if its model fails to load or its run crashes, never because of its outcome, and Appendix~\ref{app:prereg} lists every excluded cell.

\subsection{Decision models}
\label{sec:setup:decision}

A fourth backend, \texttt{decision}, speaks the System One protocol of Jev and its open-weight relatives \citep{typesafe2026choice,palmer2026kev,cloudflare2026clef}: the rendered task description and observation go out as the state, the legal actions as the criteria of a single choice question, and the returned probabilities are used exactly as the log-likelihoods of \texttt{prompt\_score}. Table~\ref{tab:systems} lists the systems and what is verified about each from its maker's cards, announcements and repositories; every performance figure quoted there is the maker's own. As with the rest of the study, the comparison includes only the systems that run on the laptop: Kev at 0.8B and 4B through its own local server, Laya through its Python package, with an adapter from its interface to the protocol, and SemIf on Qwen3.5-4B, whose readout our \texttt{score\_letter} condition also reproduces on the same checkpoint. The configuration \texttt{configs/decision\_models.yaml} evaluates them on the seeds of the main sweep, next to the Qwen3.5 0.8B and 4B base checkpoints converted by our own methods: scoring with name labels and with letter labels, the probe, and LoRA with 40 training episodes. The Qwen2.5 and Qwen3 rows of the main sweep join them at report time. Jev is offered only as a hosted service. Clef, Clef-flash and Perplexity's decider are served through hosted APIs, and their open weights, like Kev's at 9B and 27B, are larger than the 8B that the local backend holds. All of them are left to future work (Section~\ref{sec:future}), and the table marks them as such. Only scoring applies to a decision model; it cannot generate, so the free-reply and generation conditions are skipped for it.

Two rows double as validations. SemIf's readout is our letter-scoring condition, so the two should agree on the same checkpoint up to prompt wording. Kev at 0.8B and 4B sits on the same Qwen3.5 bases as our LoRA conversions and differs from them mainly in its training data, which isolates what the training target contributes, although it also reads its answer through a pointer head rather than the language-model head. Every system is reduced to the same numbers: normalized score per task with its mean and worst task, oracle agreement, calibration of the probability assigned to the chosen action against the oracle (expected calibration error and Brier score), latency per decision on the laptop (median and 95th percentile), tokens, parameters and openness.

\begin{table}[t]
\centering\small
\caption{Decision models compared with the open-weight conversions. Sizes, bases, licenses and training descriptions are from the makers' cards and announcements. ``In this study'' says how a system runs on the laptop, or marks it as future work when it needs a hosted API or more memory than the laptop has. Jev's base and training data are undisclosed.}
\label{tab:systems}
\begin{tabularx}{\linewidth}{@{}l l X X X@{}}
\toprule
System & Maker, license & Base and size & How it decides (as reported) & In this study\\
\midrule
Jev \citep{typesafe2026jev} & TypeSafe AI; closed & undisclosed & ``System One'' model; reinforcement learning for calibrated decisions on synthetic data & future work (API)\\
Kev \citep{palmer2026kev} & J.~Palmer; Apache-2.0 & Qwen3.5 0.8B/4B/9B-Base with rank-16 LoRA and a pointer head; Qwen3.8-27B full fine-tune & trained on a 10k-example decision set plus generated policy examples; System One server & 0.8B and 4B: local server; 9B and 27B: future work\\
Clef, Clef-flash \citep{cloudflare2026clef} & Cloudflare; Apache-2.0 & Qwen3.8-27B, Qwen3.5-9B, frozen & routing head with low-rank adapters; Jev-compatible & future work (API)\\
pplx-decider v1.1 \citep{perplexity2026decider} & Perplexity; Apache-2.0 reported, gated & Qwen3.8-27B fine-tune & Decisions API with its own schema & future work (API, adapter)\\
Laya \citep{convai2026laya} & Convai Innovations; Apache-2.0 & ModernBERT-large 421M; mmBERT-base 322M & encoder with a decision head; Choice, Score and Noul primitives & local package (adapter)\\
SemIf \citep{lee2026semif} & T.~Lee; MIT code & frozen Qwen3 0.6B, MiniCPM5 2B, Qwen3.5-4B & option-logit readout, no training & local, on Qwen3.5-4B; also our letter scoring\\
\bottomrule
\end{tabularx}
\end{table}

This comparison is exploratory. It was added after the first draft of the pre-registration and is kept out of its tests, and we state no hypothesis about it. The decision models were trained on data we have not inspected, so contamination cannot be ruled out, and they meet our tasks through an interface designed for classification panels. What the comparison can show is whether a conversion that anyone can reproduce on open weights reaches the decision quality of purpose-built systems of similar size on tasks with a known optimum, and at what cost on the same machine.

\subsection{Pre-registration}
\label{sec:setup:prereg}

The hypotheses, predictions and falsification criteria (Appendix~\ref{app:prereg}) are written before the main sweep and frozen, together with the prompts and the analysis code, at the git tag \texttt{prereg-v1}, which is created after the pilot on the laptop. Prompts may change only during the pilot and only against pilot seeds. Every run after the tag is confirmatory, and any departure from the frozen design is recorded with its date and reason in the deviations table of \texttt{docs/preregistration.md} and in Appendix~\ref{app:prereg}. Restricting the design to one laptop was decided before the tag and before any confirmatory run, and it is the first entry of that table. \todo{Commit hash and date of the prereg-v1 tag; deviations, if any.}
